\lecture{5.5}{Candidate Estimator Selection}
\label{lec:5-5}

\subsection{Pre-Specifying the Candidate Library}

A candidate library is a set of regression algorithms chosen before looking at which
one gives the desired causal result. For outcome regression, the library might contain
linear regression, logistic regression, generalized additive models, random forests,
lasso, regression trees, splines, and neural networks. For propensity score estimation,
the same broad idea applies, but the outcome being predicted is treatment $A$ rather
than outcome $Y$.

\begin{definition}[Candidate library]
A candidate library is a pre-specified collection
\[
\mathcal L=\{\mathcal A_1,\ldots,\mathcal A_K\}
\]
of regression algorithms. When algorithm $\mathcal A_k$ is trained on the observed
data, it produces a fitted regression estimator $\Qbar_{n,k}$ or $g_{n,k}$.
\end{definition}

The library should be broad enough to include methods with different bias-variance
profiles. A useful library often contains both simple stable methods and flexible
adaptive methods. The simple methods protect against overfitting in small samples;
the flexible methods protect against bias when the true regression is complex.

\subsection{Comparing Regressions Using Cross-Validated Risk}

For each candidate algorithm $k$, compute its cross-validated risk:
\[
R_{CV,n}(\Qbar_{n,k})
=
\frac{1}{V}\sum_{v=1}^V
\frac{1}{|\mathcal V_v|}
\sum_{i\in\mathcal V_v}
L(\Qbar_{n,k,v})(W_i,A_i,Y_i),
\]
where $\Qbar_{n,k,v}$ is algorithm $k$ fit on the training data excluding validation
fold $\mathcal V_v$.

The candidate with the smallest cross-validated risk is selected:
\[
k_{CV}
=
\argmin_{k\in\{1,\ldots,K\}}
R_{CV,n}(\Qbar_{n,k}).
\]
The selected algorithm is then refit on the full dataset to produce the final nuisance
function estimate.

\subsection{Discrete Super Learner}

\begin{definition}[Discrete Super Learner]
The discrete Super Learner is the cross-validation selector that chooses the single
candidate algorithm with the lowest cross-validated risk and then refits that algorithm
on the full sample.
\end{definition}

The word ``discrete'' means that the final estimator is one member of the candidate
library, not a weighted combination. If the library contains $K$ algorithms, the
discrete Super Learner chooses one of those $K$ algorithms.

\begin{estimationbox}
Discrete Super Learner for an outcome regression:
\begin{enumerate}
\item Pre-specify candidate algorithms $\mathcal A_1,\ldots,\mathcal A_K$.
\item Choose a valid loss function for $\Qbar$, such as squared-error loss or
negative log-likelihood loss.
\item Compute cross-validated risk for each candidate.
\item Select the candidate with the smallest cross-validated risk.
\item Refit the selected candidate on the full sample.
\item Use the resulting $\Qbar_n$ inside the causal estimator.
\end{enumerate}
\end{estimationbox}

\subsection{Avoiding Post-Hoc Model Checking}

Traditional regression workflows often involve fitting a model, checking residual
plots or diagnostic tables, adding or removing terms, and refitting. This can be useful
for scientific exploration, but it creates problems for confirmatory estimation and
inference if the final model is chosen after seeing the data.

The issue is not that diagnostics are bad. The issue is that ordinary standard errors
and confidence intervals usually treat the final model as if it had been chosen in
advance. They do not account for the uncertainty introduced by trying many models
and selecting one.

Super Learning addresses this by moving model choice into a pre-specified algorithm.
If one would consider a main-terms model, an interaction model, a spline model, and a
random forest after inspecting diagnostics, those options can be placed in the library
in advance. Cross-validated risk then supplies the selection rule.

\begin{keyidea}
The analysis plan should specify the library, loss function, cross-validation scheme,
and selection rule before examining which candidate gives the most favorable causal
estimate.
\end{keyidea}

\subsection{Using Cross-Validation for Outcome Regression and Propensity Score Estimation}

The same framework applies to the two central nuisance functions:
\[
\Qbar(a,w)=\E(Y\given A=a,W=w),
\qquad
g_1(w)=\Prob(A=1\given W=w).
\]
For outcome regression, the validation loss compares predicted outcomes with
observed outcomes. For propensity score estimation, the validation loss compares
predicted treatment probabilities with observed treatment values.

For a binary treatment, a common propensity-score loss is the Bernoulli negative
log-likelihood:
\[
L(g_n)(w,a)
=
-a\log\{g_n(w)\}-(1-a)\log\{1-g_n(w)\}.
\]
This loss is minimized by the true conditional treatment probability
$\Prob(A=1\given W=w)$.

It is common for different algorithms to perform best for different nuisance
functions. A highly flexible method might predict the outcome well but estimate the
propensity score poorly, or vice versa. Therefore, the outcome regression library and
propensity score library should be evaluated separately, with losses appropriate to
each target.

\begin{implementationbox}
Conceptual R workflow:
\begin{lstlisting}[language=R]
# Outcome regression library:
# compare algorithms by cross-validated loss for predicting Y from A and W.

# Propensity score library:
# compare algorithms by cross-validated loss for predicting A from W.

# The selected outcome algorithm and selected propensity algorithm
# need not be the same.
\end{lstlisting}
\end{implementationbox}
